\documentstyle[% 


righttag% 


,11pt% 


,amssymb% 


,russian%              remove in Eng. 


,izvru%                Set izven% in Eng. 


]{amsart} 


 


%  


 


\newcommand{\tts}[1]{} 


 


%\hoffset=-15mm%                  For Eng. 


\setcounter{page}{59} 


\god{1999} 


\nomer{10~(449)} %                        Remove in Eng. 


%\nomer{Vol.\,43, No.\,10}%               Set in Eng. 


%\str{\pageref{begin}--\pageref{end}}%  Set in Eng. 


\udk{517.518.84} 


 


\author{\it П.А.\,Шулиманов} 


% NB:   ^^ remove in Eng. 


\title{Оценка снизу приближений на полуоси целых функций  


дробно-рациональными специального вида} 


 


\date{} 


%\pagestyle{myheadings}%         Set in Eng. 


\pagestyle{plain} %              Remove in Eng. 


\begin{document} 


%\thispagestyle{plain}%          Set in Eng. 


\maketitle 


\label{begin} 


 


Пусть $f(z)=\sum\limits^\infty_{k=0}a_k z^k$ --- целая функция с 


неотрицательными 


коэффициентами $a_k$ ($k\ge1$), причем $a_0>0$. Примем 


величину 


$\big\|\frac{1}{f}-\frac{P_m}{Q_n}\big\|_{L_\infty[0,\infty)}$ за 


критерий качества приближения указанного вида функций 


дробно-рациональными функциями $P_m/Q_n$, где 


$P_m(x)=\sum\limits^m_{k=0}\ov{a_k}e^{kx}$, $0\le m\le n$, 


не убывает на $[0,+\infty)$, а $Q_n(x)=\sum\limits^n_{k=0}b_k e^{kx}$ 


($b_k\ge 0$). Характеристикой 


роста целой функции $f$ является функция 


$$ 


M(r)=M_f(r)=\max_{|z|\le r}|f(z)|,\quad r\ge 0. 


$$ 


Под порядком целой функции понимают величину 


$$


\limsup_{r\to\infty}\frac{\ln\ln M_f(r)}{\ln r} 


=\rho\quad (0\le\rho\le\infty). 


$$


Если $0<\rho<\infty$, то тип $\tau$ и нижний тип $\omega$ целой функции 


$f$ определяют как 


$$ 


\lim_{r\to\infty}\frac[0pt]{\sup}{\inf} 


\frac{\ln M_f(r)}{r^\rho}=\frac[0pt]{\tau}{\omega} 


\quad (0\le\omega\le\tau\le\infty). 


$$ 


 


Через $H^+(\rho,\omega,\tau)$ обозначим класс целых функций $f$ 


заданного порядка 


$\rho$ ($0<\rho<\infty$), типа $\tau$ и нижнего типа $\omega$ 


($0<\omega\le\tau<\infty$)  


с неотрицательными коэффициентами Тейлора $a_k$ ($k\ge 1$), $a_0>0$. 


В данной заметке получена новая оценка снизу величины  


$\big\|\frac{1}{f}-\frac{P_m}{Q_n}\big\|_{L_\infty[0,\infty)}$ для 


функций $f\in H^+(\rho,\omega,\tau)$, $\rho>1$, которая как следствие 


усиливает результат 


работы [1]. Известна  


 


\begin{thm}[{\series{m}\shape{n}\selectfont\cite{1}}] Пусть 


$$ 


f(z)=\sum^\infty_{k=0}a_k z^k,\quad a_0>0,\quad 


a_k\ge 0\quad (k\ge 1) 


$$ 


--- целая функция порядка $\rho=2$, типа $\tau$ и нижнего типа $\omega$ 


$(\frac{1}{25}\le\omega\le\tau<\infty)$ 


или порядка $\rho$ $(2<\rho<\infty)$, типа $\tau$ и нижнего типа 


$\omega$ $(0<\omega\le\tau<\infty)$. Тогда невозможно найти полиномы 


вида 


$\sum\limits^n_{k=0}b_k e^{kx}$, $b_k\ge 0$, для которых 


$$ 


\liminf_{n\to\infty}\bigg\{\bigg\|\frac{1}{f}- 


\frac{1}{\sum\limits^n_{k=0}b_k e^{kx}} 


\bigg\|_{L_\infty[0,\infty)}\bigg\}^{\frac{\rho\omega}{n^2\tau}} 


\le e^{-1}. 


$$ 


\end{thm} 


 


{\bf Замечание.} На самом деле приведенная в теореме оценка для  


функций порядка $\rho=2$, типа $\tau$ и нижнего типа $\omega$ 


($\frac{1}{25}\le\omega\le\tau<\infty$) 


доказана в [1] при более сильных ограничениях на 


параметры. А именно, в заключительной 


части доказательства неравенства (13) в [1] появляется ограничение, 


которое после несложных преобразований приводит к 


ограничению $\big(\frac{2}{\sqrt{\omega}}-\frac{19}{21\sqrt{2}}\big) 


\sqrt{\tau}<1$ на параметры $\omega$ и $\tau$. 


 


В дальнейшем будем использовать обозначения 


\begin{equation} 


\varphi_\varepsilon=\exp\bigg\{d_\varepsilon(\delta) 


\frac{\delta-1}{d_\varepsilon(\delta)-1}\bigg\{ 


\frac{\delta-1}{(d_\varepsilon(\delta)-1)\omega(1-\varepsilon)} 


\bigg\}^{\frac{1}{\rho-1}}\bigg\}, 


\end{equation} 


где 


\begin{gather} 


d_\varepsilon(\delta)= 


\frac{\delta^\rho\omega(1-\varepsilon)}{\tau(1+\varepsilon)}, 


\quad 0\le\varepsilon<1,\notag \\ 


c=c(\varepsilon,\omega,\tau,\rho)=\min\bigg\{ 


\varphi_\varepsilon(\delta):d_\varepsilon(\delta)= 


\frac{\delta^\rho\omega(1-\varepsilon)}{\tau(1+\varepsilon)} 


>1\bigg\}=\exp\bigg\{ 


\frac{(2\delta_0-1)^\rho}{\rho^\rho\omega(1-\varepsilon) 


d_\varepsilon(\delta_0)}\bigg\}^{\frac{1}{\rho-1}} 


\end{gather} 


и $\delta_0$ --- корень уравнения 


\begin{equation} 


(2\delta-1)(d_\varepsilon(\delta)-1)= 


\rho d_\varepsilon(\delta)(\delta-1). 


\end{equation} 


 


Основным результатом работы является  


 


\begin{thm} 


Пусть $0<\varepsilon_0<1$, $0\le\varepsilon<1$ и $f\in 


H^+(\rho,\omega,\tau)$, $\rho>1$. 


Тогда существует $n_0=n_0(\varepsilon_0,\varepsilon)$ такое, что для 


всех $n\ge n_0$ и 


неубывающих на $[0,\infty)$ полиномов вида 


$P_m(x)=\sum\limits^m_{k=0}\ov{a_k}e^{kx}$, $0\le m\le n$, и 


$Q_n(x)=\sum\limits^n_{k=0}b_k e^{kx}$ 


$(b_k\ge0)$ справедливо неравенство 


\begin{equation} 


\bigg\|\frac{1}{f}-\frac{P_m}{Q_n}\bigg\|_{L_\infty[0,\infty)} 


\ge c_1(1-\varepsilon_0)c^{-n^{\frac{\rho}{\rho-1}}}, 


\end{equation} 


где величина $c=c(\varepsilon,\omega,\tau,\rho)$ определена формулой 


$(2)$, а константа 


$c_1=c_1(\varepsilon,\omega,\tau,\rho)$ зависит только от указанных 


величин. 


\end{thm} 


 


{\bf Следствие.} В условиях теоремы 2 справедливо неравенство 


$$ 


\liminf_{n\to\infty}\bigg\{ 


\bigg\|\frac{1}{f}-\frac{P_m}{Q_n}\bigg\|_{L_\infty[0,\infty)} 


\bigg\}^{\frac{1}{n^{\rho/(\rho-1)}}}\ge c^{-1}_2, 


$$ 


где $c_2=c(0,\omega,\tau,\rho)$ определяется формулой (2) при 


$\varepsilon=0$. 


\medskip 


 


В частном случае при $m=0$ ($\rho=2$) получаем усиление по порядку 


при $n\to\infty$ оценки А.Р.\,Редди [1] и расширение ее на более  


широкий класс параметров. Действительно, сравним полученный в  


следствии результат с теоремой 1 для функций $f\in H^+(2,\omega,\tau)$. 


Для этого положим в (1) $\rho=2$, $\varepsilon=0$. Тогда имеем 


$$ 


\varphi_0(\delta)=\exp\bigg\{\frac{\tau}{\omega^2}\bigg( 


\frac{\delta^2-\delta}{\delta^2-\tau/\omega}\bigg)^2\bigg\} 


$$ 


при значениях $\delta>\sqrt{\frac{\tau}{\omega}}$. Покажем, что она 


меньше соответствующей 


оценки $\exp\big\{\frac{\tau}{2\omega}\big\}$ из теоремы~1 для всех 


значений $\delta\ge\frac{\tau}{\omega}$, где  


$2.44\dotsc\le\omega\le\tau<\infty$ (см. замечание к теореме 1). 


Логарифмируя  


указанные оценки, получаем неравенство 


$$ 


\ln\varphi_0(\delta)=\frac{\tau}{\omega^2}\bigg( 


\frac{\delta^2-\delta}{\delta^2-\tau/\omega}\bigg)^2< 


\frac{\tau}{2\omega}, 


$$ 


которое выполняется, т.к. 


$\frac{\delta^2-\delta}{\delta^2-\tau/\omega}\le 


1<\sqrt{\frac{\omega}{2}}$. Ясно, что 


$$ 


\min_{\delta\ge\sqrt{\frac{\tau}{\omega}}} 


\ln\varphi_0(\delta)=\ln\varphi_0(\delta_0)< 


\frac{\tau}{2\omega}, 


$$ 


где величина $\delta_0=\frac{\tau}{\omega}+\big( 


\frac{\tau^2}{\omega^2}-\frac{\tau}{\omega}\big)^{1/2}$. Отметим, что 


для функций $f\in H^+(2,\frac{1}{4},\frac{1}{4})$ 


полученный в следствии порядок оценки снизу является  


логарифмически точным, т.к. из следствия имеем 


$$ 


\liminf_{n\to\infty}\bigg\{ 


\bigg\|\frac{1}{f}-\frac{P_m}{Q_n}\bigg\|_{L_\infty[0,\infty)} 


\bigg\}^{1/n^2}\ge\frac{1}{e}, 


$$ 


а в ([1], с.\,254) приведен пример функции из этого класса, для  


которой выполняется неравенство 


$$ 


\limsup_{n\to\infty}\bigg\{ 


\bigg\|\frac{1}{S_n}-\frac{1}{f}\bigg\|_{L_\infty[0,\infty)} 


\bigg\}^{1/n^2}\le e^{-1/2},\quad 


\text{где} 


\quad S_n(z)=\sum^n_{k=0}a_k e^{kz}. 


$$ 


Отметим, что при $m=0$ ($\rho>0$) следствие усиливает логарифмический 


порядок при $n\to\infty$ соответствующей оценки из теоремы~1. 


 


{\bf Доказательство теоремы 2.} Пусть 


\begin{equation} 


A=\frac{1}{(\omega(1-\varepsilon))^{1/(\rho-1)}}\bigg\{ 


\frac{\delta_0-1}{d_\varepsilon(\delta_0)-1} 


\bigg\}^{\frac{\rho}{\rho-1}},\quad 


B=d_\varepsilon(\delta_0)A,\quad 


\sigma=\{d_\varepsilon(\delta_0) 


\}^{\frac{d_\varepsilon(\delta_0)}{d_\varepsilon(\delta_0)-1}}, 


\end{equation}  


где $\delta_0$ --- корень уравнения (3). Доказательство будем вести от 


противного. 


Предположим, что неравенство в (4) при некотором $\varepsilon_0>0$ 


неверно. Тогда существуют бесконечная последовательность целых  


чисел $1\le n_1<n_2<\dotsb$ и последовательность ``полиномов'' 


$P_{n_k}$ и $Q_{n_k}$ 


указанного вида таких, что выполняется неравенство 


\begin{equation} 


\bigg\|\frac{1}{f}-\frac{P_{n_k}}{Q_{n_k}}\bigg\|_{L_\infty[0,\infty)} 


<c_1(1-\varepsilon_0)\exp\{-B n_k^{\frac{\rho}{\rho-1}}\}< 


\exp\{-B n_k^{\frac{\rho}{\rho-1}}\}. 


\end{equation} 


Поскольку $f$ --- целая функция и $a_k\ge 0$ ($k\ge0$), то 


$\lim\limits_{x\to\infty}f(x)=\infty$ и для  


каждого $n_k$ существует число $r_k$ со свойством 


\begin{equation} 


f(r_k)=\sigma^{\frac{1}{d_\varepsilon(\delta_0)}}\exp 


\{A n_k^{\frac{\rho}{\rho-1}}\}, 


\end{equation} 


где величины $A$ и $\sigma$ определены в формуле (5). Учитывая 


определение 


нижнего типа $\omega$, из условия (7) получаем неравенство 


\begin{multline} 


r_k\le\bigg\{\frac{A}{\omega(1-\varepsilon)}\bigg\} 


^{1/\rho}n_k^{\frac{1}{\rho-1}}\bigg\{1+ 


\frac{\ln\sigma}{A d_\varepsilon(\delta_0) 


n^{\frac{\rho}{\rho-1}}}\bigg\}^{1/\rho}\le \\ 


% 


\le\bigg\{\frac{A}{\omega(1-\varepsilon)}\bigg\}^{1/\rho} 


n_k^{\frac{1}{\rho-1}}\bigg\{ 


1+\frac{\ln\sigma}{A d_\varepsilon(\delta_0) 


n_k^{\frac{\rho}{\rho-1}}}\bigg\}=\bigg\{ 


\frac{A}{\omega(1-\varepsilon)}\bigg\}^{1/\rho} 


n_k^{\frac{1}{\rho-1}}+\ln\sigma_1^{\frac{1}{n_k}}, 


\end{multline} 


где $\sigma_1=\sigma^{1/t}$, $t=(\omega(1-\varepsilon))^{1/\rho} 


A^{\frac{\rho-1}{\rho}}d_\varepsilon(\delta_0)$. Из неравенства (6) 


легко следует 


\begin{equation} 


\bigg|\frac{Q_{n_k}(r_k)}{P_{n_k}(r_k)}\bigg|< 


\sigma^{\frac{1}{d_\varepsilon(\delta_0)}}\exp 


\{A n_k^{\frac{\rho}{\rho-1}}\}\bigg\{1- 


\frac{\sigma^{\frac{1}{d_\varepsilon(\delta_0)}}} 


{\exp\{(B-A)n_k^{\frac{\rho}{\rho-1}}\}}\bigg\}^{-1}. 


\end{equation} 


Положим $x_k=r_k\delta_0$, где числа $r_k$ определены в (7), а 


$\delta_0$ --- корень  


уравнения (3). В данном случае $M(r)=f(r)$, поэтому из определения 


типа $\tau$, нижнего типа $\omega$ и из соотношения (7) для достаточно  


больших значений $k$ следует, что 


\begin{equation} 


f(x_k)=f(r_k\delta_0)\ge\{f(r_k)\}^{d_\varepsilon(\delta_0)}= 


\sigma\exp\{B n_k^{\frac{\rho}{\rho-1}}\}. 


\end{equation} 


Используя неравенства (8) и (9), получаем оценку отношения 


$\frac{Q_{n_k}}{P_{n_k}}$ в  


точках $x_k=r_k\delta_0$, а именно, 


$$


\bigg|\frac{Q_{n_k}(r_k\delta_0)}{P_{n_k}(r_k\delta_0)}\bigg|\le 


\exp\{n_k(\delta_0-1)r_k\}\bigg| 


\frac{Q_{n_k}(r_k)}{P_{n_k}(r_k)}\bigg|\le \\ 


% 


\sigma_1^{\frac{1}{n_k}} 


\sigma^{\frac{1}{d_\varepsilon(\delta_0)}}\exp 


\{B n_k^{\frac{\rho}{\rho-1}}\}\bigg\{1- 


\frac{\sigma^{\frac{1}{d_\varepsilon(\delta_0)}}} 


{\exp\{(B-A)n_k^{\frac{\rho}{\rho-1}}\}}\bigg\}^{-1}. 


$$


 


Покажем теперь, что полученное неравенство и (10) противоречат  


предположению (6). Действительно, неравенство 


\begin{multline*} 


\bigg|\frac{P_{n_k}(r_k\delta_0)}{Q_{n_k}(r_k\delta_0)}- 


\frac{1}{f(r_k\delta_0)}\bigg|\ge\bigg| 


\frac{P_{n_k}(r_k\delta_0)}{Q_{n_k}(r_k\delta_0)}\bigg|- 


\frac{1}{f(r_k\delta_0)}\ge \\ 


% 


\ge\frac{1-\sigma^{\frac{1}{d_\varepsilon(\delta_0)}}\exp\{ 


-(B-A)n_k^{\frac{\rho}{\rho-1}}\}} 


{\sigma_1^{\frac{1}{n_k}} 


\sigma^{\frac{1}{d_\varepsilon(\delta_0)}}\exp 


\{B n_k^{\frac{\rho}{\rho-1}}\}}- 


\frac{1}{\sigma\exp\{B n_k^{\frac{\rho}{\rho-1}}\}}\ge 


(1-\varepsilon_0)c_1\exp\{-B 


n_k^{\frac{\rho}{\rho-1}}\} 


\end{multline*} 


справедливо для всех значений $k\ge k_0(\varepsilon_0)$. Это 


противоречит левому  


из неравенств (6). Следовательно, справедливо неравенство (4).  


\medskip 


 


{\bf Замечание.} Отметим, что результат данной теоремы был анонсирован 


автором в [2]. 


 


\begin{thebibliography}{9} 


 


\bibitem{1} 


Reddy A.R. {\em Some results in Chebyshev rational approximation} // 


J. Approxim. Theory. -- 1976. -- V.\,18. -- \No\,3. -- Р.\,251--263. 


\bibitem{2} 


Шулиманов П.А. {\em Оценки снизу приближений на полуоси целых  


функций дроб\-но-рациональными специального вида} // Республиканск. 


конф. по экстремальным задачам теории приближений 


и их прилож. Тез. докл. -- Киев, 1990. -- 


С.\,147. 


 


\end{thebibliography} 


 


\vskip 2cm 


 


\post{\it Уральская государственная}{\it Поступила} 


\post{\it сельскохозяйственная академия}{09.01.1996} 


 


\label{end} 


\end{document} 


